% Quarantine-only proof.
% Date: 29 August 2026.
% This file is not a live-book input and makes no novelty or priority claim.

\begin{proof}[Proof of
\cref{lem:p213-trivial-gauss-rigid-zero-atom}]
The normal locus is Zariski open by \cite[\S11.3.1]{CLD25}.  After
shrinking $\Omega$ about $y$, we may therefore assume that
\[
 \Omega\subseteq\operatorname{Int}(X/K),
 \qquad
 \Omega\ \text{is normal and purely one-dimensional}.
\]
An open analytic subspace of the analytic interior is boundaryless by
the boundary-composition formula \cite[(11.4.1.3)]{CLD25}.

Fix $0<r<1$.  Let
\[
 K_r=K((\varpi)),\qquad |\varpi|=r,
\]
with its $\varpi$-adic Gauss norm.  In the notation of
\cite[Example~11.1.3]{CLD25}, this is the field attached to the one-term
polyradius $(r)$.  Since $K$ is trivially valued and $r\ne1$, this
polyradius is $K$-free.  Thus $K_r/K$ is analytically separable by
\cite[\S11.1.9]{CLD25}.  It is also universally multiplicative, and hence
the projection
\[
 \pi:\Omega_r:=\Omega\widehat\otimes_K K_r\longrightarrow\Omega
\]
has the continuous Shilov section
$\sigma:\Omega\to\Omega_r$ described in
\cite[\S12.7.4--Corollary~12.7.5]{CLD25}.

We record all geometric properties needed below.  Analytic separability,
the scalar-extension statements in \cite[\S\S11.1.9 and 11.3.1]{CLD25},
and the $R_1+S_2$ criterion show that $\Omega_r$ is normal.  Every
irreducible component after scalar extension has the same dimension as,
and maps onto, an irreducible component before scalar extension by
\cite[\S11.2.5]{CLD25}; hence $\Omega_r$ is purely one-dimensional.
Boundarylessness is stable under analytic base change by
\cite[\S11.4.1]{CLD25}.  Goodness and separatedness are preserved by
scalar extension.  Strictness follows directly on affinoid charts: the
base change of a quotient of
$K\langle T_1,\ldots,T_m\rangle$ is a quotient of
$K_r\langle T_1,\ldots,T_m\rangle$.

The lift of $y$ is rigid without any separability hypothesis on its
residue field.  Indeed, put $L=\mathscr H(y)$.  Rigidity says that $L/K$
is finite, and the valuation on $L$ is trivial.  Coefficientwise
expansion gives an isometric isomorphism of Banach $K_r$-algebras
\begin{equation}
\label{eq:p213-trivial-gauss-residue-tensor}
             L\widehat\otimes_K K_r\simeq L((\varpi)).
\end{equation}
To see the isometry explicitly, choose a $K$-basis of $L$.  It is
orthonormal for the trivial norm, and the Gauss norm of a Laurent series
with coefficients in $L$ is the maximum of the Gauss norms of its finitely
many coordinate series.  The right-hand side of
\eqref{eq:p213-trivial-gauss-residue-tensor} is a field finite over $K_r$.
Consequently the scalar fibre above $y$ is a singleton, namely
$\sigma(y)$, and
\[
             \mathscr H(\sigma(y))\simeq L((\varpi)).
\]
Thus $\sigma(y)$ is a rigid interior point of the good separated strict
normal pure one-dimensional $K_r$-analytic space $\Omega_r$.  This
argument is unchanged when $L/K$ is purely inseparable.

Set $u_r=\pi^*u$.  It is continuous and CLD-psh by
\cite[Proposition~8.1.5]{CLD25}.  Put
\[
                         \mu_r=\mathrm d'\mathrm d''[u_r].
\]
The scalar-extension support statement
\cite[Proposition~16.8.2]{CLD25}, applied with $n=p=1$ and the smooth
$(0,0)$-form $1$, gives
\begin{equation}
\label{eq:p213-trivial-gauss-shilov-support}
                         \operatorname{supp}(\mu_r)
                         \subseteq\sigma(\Omega).
\end{equation}
The pushforward identity recalled in the same proposition from
\cite[Proposition~8.7.5]{CLD25} gives
\begin{equation}
\label{eq:p213-trivial-gauss-pushforward}
                         \pi_*\mu_r=\mu_u.
\end{equation}
Both are positive Radon measures; positivity and the Radon assertion also
follow directly from \cite[Lemma~8.5.5 and Proposition~7.2.8]{CLD25}.
Since $\pi\circ\sigma$ is the identity,
\eqref{eq:p213-trivial-gauss-shilov-support} and
\eqref{eq:p213-trivial-gauss-pushforward} imply
\begin{equation}
\label{eq:p213-trivial-gauss-atom-transfer}
             \mu_u(\{y\})
             =\mu_r\bigl(\pi^{-1}(\{y\})\bigr)
             =\mu_r(\{\sigma(y)\}).
\end{equation}
The nontrivially valued rigid-point estimate
\cref{lem:p21p3-componentwise-normal-rigid-zero-atom} applies to
$u_r$ at $\sigma(y)$ and makes the last term zero.  This proves
\eqref{eq:p213-trivial-gauss-rigid-zero-atom}.
\end{proof}
